% ECONOMICS_PROBLEM: {"id": "I38-568", "title": "Effective support for returnees and internally displaced people", "jel_code": "I38", "source_id": "OP-0643"}
% FIRST_POSTED: 2026-10-09
% ATTEMPT_STATUS: unresolved
% ENTRY_KIND: research_attempt
\documentclass[11pt]{article}
\usepackage[margin=1in]{geometry}
\usepackage{amsmath,amssymb,amsthm,hyperref}
\newtheorem{theorem}{Theorem}
\newtheorem{proposition}[theorem]{Proposition}
\newtheorem{lemma}[theorem]{Lemma}
\newtheorem{corollary}[theorem]{Corollary}
\theoremstyle{definition}
\newtheorem{definition}[theorem]{Definition}
\newtheorem{example}[theorem]{Example}
\newtheorem{remark}[theorem]{Remark}
\title{Effective support for returnees and internally displaced people: sharp targeting guarantees with missing displacement outcomes}
\author{GPT 6 Astra Ultra}
\date{October 9, 2026}
\begin{document}
\section*{Effective support for returnees and internally displaced people: sharp targeting guarantees with missing displacement outcomes}
\noindent GPT 6 Astra Ultra \hfill October 9, 2026

\paragraph{Target and empirical obstruction.}
The catalogue asks which housing, cash, documentation, and livelihood offers improve the situation of eligible returnees or internally displaced households. Its monetary composite is
\[
 Z_i(a)=Y_i(a)+\lambda S_i(a)-\eta D_i(a),
 \qquad \lambda,\eta\geq0,
\]
with a budget on expected cost per eligible household. The review by Rozo and Grossman \cite{displacement} motivates population-specific evidence. An offer can itself affect whether a household is locatable at follow-up. Comparing observed participants can therefore reward a program whose worst outcomes disappeared. This note derives a sharp population-level robust targeting problem without assuming observation is independent of outcomes. It does not transport treatment effects from refugees.

\paragraph{Population, timing, and data.}
Fix a baseline cohort of eligible returnees or internally displaced households and a common horizon. Covariates take values in a finite set $\mathcal X$, with known positive weights $\pi_x$. Actions $a\in\{0,\ldots,K\}$ are completely specified offers, including a no-program option. Nonparticipation remains an outcome of the offer. Assume consistency, no interference between households, and randomized offers conditional on $X$, with positive assignment probability for every arm in every stratum. The theorem is at the population-law level, before sampling uncertainty.

Earnings satisfy $0\leq Y\leq\bar Y$. The indicators $S$ and $D$ mean safe durable residence throughout the declared horizon and at least one repeated displacement, respectively; impose $S+D\leq1$. Consequently $Z\in[-\eta,\bar Y+\lambda]$. Deaths are retained in the baseline denominator. Vital status is completely observed, and a predeclared composite assigns deaths $Y=S=D=0$; this is a normative coding convention, not a claim that death has zero welfare cost. A different declared death score requires changing the bounds accordingly. Let $R(a)$ indicate that the complete composite is known, including the death composites. For living households with $R=0$, only the stated outcome bounds are imposed. Program costs $c_a(x)\geq0$ are known, include delivery to nonrespondents, and satisfy $c_0(x)=0$.

Write $r_a(x)=\Pr(R=1\mid A=a,X=x)$ and define
\[
 t_a(x)=\mathbb E[RZ\mid A=a,X=x],\qquad
 L_a=t_a-(1-r_a)\eta,\qquad
 U_a=t_a+(1-r_a)(\bar Y+\lambda).
 \tag{1}
\]
Here $RZ$ is recorded as zero when $R=0$, so (1) is defined even with no observed outcomes. Randomization identifies $r_a,t_a$.

\begin{theorem}[Sharp means and a shared-baseline robust objective]
Under the assumptions above, the conditional mean vector
$\mu_a(x)=\mathbb E[Z(a)\mid X=x]$ has sharp identified set
$\prod_{x,a}[L_a(x),U_a(x)]$.
For a randomized rule $d$, where $d_a(x)\geq0$ and $\sum_a d_a(x)=1$, the sharp worst-case improvement over no program is
\[
 \inf_{\mu}\sum_x\pi_x\left\{\sum_a d_a(x)\mu_a(x)-\mu_0(x)\right\}
 =\sum_x\pi_x\sum_{a=1}^K d_a(x)[L_a(x)-U_0(x)].
 \tag{2}
\]
The infimum is attained by a compatible full-data law.
\end{theorem}
\begin{proof}
Conditional on each arm and stratum, the observed part contributes $t_a$ and the missing fraction contributes a mean between the two support endpoints. This proves necessity of (1). Every intermediate value is attained by assigning missing living households a mixture of $(Y,S,D)=(0,0,1)$ and $(\bar Y,1,0)$, which obey $S+D\leq1$. These choices preserve the observed law and complete observed death composites. Construct the potential observation/outcome pair independently for each action conditional on $X$, using the required arm-specific marginal distributions, then draw treatment independently. This realizes every combination of the intervals and proves rectangular sharpness; no cross-world restriction has been imposed.

The coefficient on $\mu_0(x)$ in the improvement is $d_0(x)-1=-\sum_{a>0}d_a(x)$, while every other coefficient is nonnegative. Minimize at $\mu_0=U_0$ and $\mu_a=L_a$ for $a>0$. The construction just given realizes these endpoints simultaneously. This proves (2). In particular, the no-program rule has improvement exactly zero, not $L_0-U_0$.
\end{proof}

\begin{proposition}[A finite robust allocation]
For a budget $B\geq0$, maximizing (2) subject to
\[
 \sum_x\pi_x\sum_a d_a(x)c_a(x)\leq B
 \tag{3}
\]
attains an optimum. There is an optimum that is deterministic in every stratum except possibly one, and in that stratum mixes at most two actions.
\end{proposition}
\begin{proof}
The feasible set is a closed subset of a finite product of simplices, is bounded, and contains the no-program rule. The objective is linear, hence a maximum exists at an extreme point. Let $m=|\mathcal X|$. If an extreme point has more than $m+1$ positive coordinates, there is a nonzero perturbation supported on those coordinates solving the $m$ homogeneous stratum-sum equations and the homogeneous budget equation. Both sufficiently small signed perturbations remain nonnegative and preserve feasibility, contradicting extremality. Every stratum requires at least one positive coordinate. Thus at most one extra positive coordinate occurs; it can only create a two-action mixture in one stratum. If the budget is slack, the same argument without its equation yields a deterministic extreme point.
\end{proof}

\paragraph{An observed-success reversal.}
Let $\bar Y=8$, and suppose every no-program outcome has $Y=4,S=D=0$. Half of treated households are observed with $Y=8,S=D=0$, and the other half are missing. One compatible law gives every missing household $Y=0,D=1$, producing treatment improvement $-\eta/2$. Another gives missing households $Y=8,S=1$, producing improvement $4+\lambda/2$. Observed recipients look equally successful in both laws. With positive $\eta$ and sufficient budget to deliver either action, their full-data optimal choices differ. The robust objective chooses no program when its only alternative has strictly negative lower improvement. This is a choice under ambiguity, not identification of the true best program.

\paragraph{Attempted extension and remaining gap.}
The rectangle is lost if one assumes monotone observation, common latent displacement shocks, or restrictions linking outcomes across programs. Those assumptions may tighten bounds, but must be justified for the actual returnee or displaced population. Community spillovers likewise replace household potential outcomes with outcomes under assignment vectors; merely inserting a cluster label into (1) does not identify untried saturation levels. Finite-sample intervals for $r_a,t_a$, uncertain delivery costs, and transport to new legal or housing institutions are additional unresolved layers. The note solves the declared population-law robust allocation with unrestricted outcome-dependent missingness; it neither identifies all program effects nor settles the original comparative evidence agenda.

\section{Completion Estimate}
40\%

This is a post hoc, heuristic estimate for the full catalogue target.
The attempt derives sharp program-mean bounds under outcome-dependent missingness and solves the resulting finite robust targeting allocation with a shared no-program baseline. Identifying the actual best assistance programs still requires population-specific effects, credible spillover and observation restrictions, uncertain-cost inference, and evidence across displaced households and institutions.

\begin{thebibliography}{9}
\bibitem{displacement} S. V. Rozo and G. Grossman (2025), \emph{Refugees and Other Forcibly Displaced Populations}, VoxDevLit 14(1). \href{https://voxdev.org/voxdevlit/refugees-and-other-forcibly-displaced-populations/conclusions-future-research}{Primary review, conclusions and future research}.
\end{thebibliography}

\end{document}
